\section{Scope and Methodology}\label{sec:scope}

\subsection{Scope}

This review evaluated the QuarkyLab / NetFRAME home lab across 21 assessment
dimensions: reliability, availability, recoverability, security, network
correctness, configuration consistency, monitoring quality, alert quality, storage
integrity, backup quality, power resilience, energy efficiency, capacity planning,
operational maturity, documentation quality, automation safety, AI-system
reliability, scientific-workload support, student-environment isolation, recruiter
and hiring-manager presentation, and Fortune-100 enterprise alignment.

In scope were only systems, repositories, and services that belong to the home lab:
the 7-node km-cluster plus the standalone pve1 host, Randy and the DS4246 shelf, the
Juniper and UniFi switching, OPNsense, Headscale, the RKE2 cluster, the Docker/LXC
service estate, the Ollama/Open WebUI/llm\_router AI stack, the Grafana and Wazuh
observability planes, the UPS/PDU/power layer, and the operator's GitHub
repositories. No unrelated external system was probed. External services already
used by the environment (package repositories, GitHub, container registries,
certificate authorities, DNS resolvers) were treated as passive dependencies, not
targets.

\subsection{Operating posture: read-only by classification}

Every command was classified before execution as read-only/low-risk,
read-only/resource-intensive, configuration-changing, service-affecting,
destructive, or external-target interaction. Only read-only/low-risk commands were
run without approval. No configuration-changing, service-affecting, destructive, or
external-interaction command was executed. Concretely, this review performed no
restarts, reboots, reloads, migrations, deletions, rotations, or deployments; used
no discovered credential; injected no failures; and did not enforce the VLAN 1
egress lockdown, in line with the standing decision that its prerequisites are not
yet met. Where confidence that an action was non-disruptive fell below the required
threshold, the action was not taken and the limitation was recorded (see
Section~\ref{sec:limits}).

\subsection{Evidence approach}

The review compared five state layers where possible: declared (the review
specification and prior intent), configured (files and definitions), deployed
(running processes and guests), observed (live telemetry), and documented (the
repository). Documentation was not treated as true merely because it exists, and
live output was not treated as authoritative where it represents only transient
runtime state.

A key input was the operator's own prior audit (commit \code{65f3681}, 2026-07-22),
which had already captured per-node live evidence and built a canonical fact ledger.
That corpus is one day old and was used as generated evidence to be re-verified, not
trusted blindly. This review independently re-ran the highest-stakes checks against
current live hardware.

\subsection{Independent validation (multi-method)}

For every high-severity conclusion, at least two independent read-only methods were
used. Examples actually applied in this review:

\begin{itemize}
  \item \textbf{GPU assignment:} live \code{nvidia-smi} on both R730s (E-02, E-03)
  versus the review specification and the fact ledger. The live reading settled the
  contradiction in favor of the repository (DR-01).
  \item \textbf{Storage health:} live \code{zpool status} on Randy (E-04) versus the
  ledger's canonical geometry and the prior audit's finding on \code{mpathv}.
  \item \textbf{Cluster and guests:} live \code{pvecm status} plus
  \code{pvesh get /cluster/resources} (E-01, E-06) versus the documented node and
  service placement, which surfaced the Grafana/Headscale relocation drift.
  \item \textbf{Power:} live \code{upsc} readings (E-07) versus the reliability
  assessment's UPS-imbalance finding, which confirmed UPS-A is the loaded unit.
  \item \textbf{Backups:} live PBS snapshot listing dated 2026-07-23 (E-08) versus
  the documented backup schedule.
\end{itemize}

\subsection{What was not done}

No failure injection, no restore into production, no restarts, and no live query of
OPNsense (E-23). The EX3400 exclusion was lifted mid-review when the operator supplied
credentials, so switch findings are live-verified (E-22, E-25).

Two configuration changes were made, both \emph{only} after explicit operator approval and
both verified afterwards: the ZFS device-level metric and alert rule, and the EX3400
\code{bridge-priority}. Everything else remained read-only. Findings that still rest on
documentation are marked medium-confidence and carry a precise follow-up request.

\subsection{Deliverable structure}

The deliverables are this LaTeX report (sections plus appendices), a set of
machine-readable registers under \code{data/} (evidence, assets, services, drift,
risk, recommendations, backup matrix, alert matrix), and text-based topology
diagrams under \code{diagrams/} (physical, logical, service-dependency,
recovery-dependency, power). The \code{evidence/} directory describes evidence
locations and collection methods rather than copying sensitive output.
